\documentclass[11pt]{article}
\usepackage{fontspec,amsmath,amssymb,amsthm,hyperref}
\usepackage[margin=1in]{geometry}
\setlength{\emergencystretch}{3em}
\title{Exact average-smoothness Gibbs score and Laplacian moments}
\author{LangevinLean mathematical exposition}
\date{}
\begin{document}\maketitle
\noindent Source module: \texttt{AverageLMC/ScoreMoment.lean}. Compile this proof document with XeLaTeX or LuaLaTeX. The accompanying HTML pairs the mathematical argument with the exact Lean statements and records the scope of the formalization.


The first-order part proves E\_pi norm(gradient f) squared at most p M\_av without assuming any Hessian or integration-by-parts identity. The second part proves Lemma 7 for C2 convex potentials with finite Gibbs normalization: the true score energy equals the L1 Laplacian integral, bounded by the exact Laplacian supremum and then the exact average coordinate smoothness. Strong convexity supplies normalization automatically in the main application.


\section{Obtain the exact coordinate score MGF by translation}


The coordinate quadratic model implies

\[e^{-f(x)}e^{-a\partial_j f(x)}\le e^{M_j a^2/2}e^{-f(x+ae_j)}.\]

Integrate the nonnegative functions against Lebesgue measure. Translation invariance leaves the partition function unchanged, and finite positive normalization cancels it. Replacing a by minus t produces the Gaussian MGF bound for the actual jth score coordinate, with proxy M\_j. The same finite exponential integral bound supplies exponential integrability, so this is an actual sub-Gaussian statement rather than a formal integral comparison.



\noindent Lean declarations: \texttt{coordinate\_score\_lmgf\_le}, \texttt{coordinate\_score\_subgaussian}.

\section{Sum sharp second moments without second differentiability}


Strong convexity implies m at most every coordinate Taylor supremum, so each M\_j is nonnegative. Apply the established sub-Gaussian second-moment theorem to the coordinate score: its uncentered second moment is at most M\_j, and square integrability is derived at the same time. The squared Euclidean norm is the sum of the squared coordinates. Integral linearity and the finite sum therefore give

\[\mathbb E_\pi\|\nabla f\|^2\le\sum_j M_j=pM_{\rm av}.\]

Only differentiability, a Lipschitz gradient, and positive strong convexity are used in this sharper score bound.



\noindent Lean declarations: \texttt{integral\_norm\_sq\_gradient\_gibbs\_le}.

\section{Identify coordinate derivatives and control all weighted integrals}


A normalized Gibbs integral is integrable exactly when its exponential-weighted Lebesgue integrand is integrable, since the partition function is finite and positive. The coordinate score equals the directional Frechet derivative in e\_j. Differentiating that expression identifies its derivative in e\_j with the actual diagonal Hessian entry. C2 regularity makes this coordinate score C1. Projecting the global gradient Lipschitz inequality bounds the coordinate score with the same Lipschitz constant, and hence bounds its derivative. These facts justify the weighted integration by parts used next.



\noindent Lean declarations: \texttt{integrable\_gibbs\_weighted}, \texttt{fderiv\_coordinate\_score}, \texttt{coordinate\_score\_contDiff}, \texttt{coordinate\_score\_lipschitz}.

\section{Prove the score-Hessian identity by global integration by parts}


Set psi equal to the jth score coordinate. Its first and second moments are finite by the first-order Gibbs score theorem. Its directional derivative is bounded by M, so that derivative is Gibbs integrable as well. The three unnormalized functions needed for global integration by parts are therefore integrable: the derivative of psi times e\textasciicircum{}\{-f\}, psi times the derivative of e\textasciicircum{}\{-f\}, and psi e\textasciicircum{}\{-f\}. The chain rule gives D\_j e\textasciicircum{}\{-f\}=-e\textasciicircum{}\{-f\} psi. Integration by parts cancels its minus sign and proves

\[\int(\partial_j f)^2\,d\pi=\int\partial_{jj}f\,d\pi.\]

No decay or boundary term is supplied as a hypothesis. Summing the finite coordinate identities and using trace(Hessian f)=Laplacian f proves the full score-Laplacian identity. The partition-based statement works whenever normalization is finite; the strongly convex wrapper discharges it internally.



\noindent Lean declarations: \texttt{integral\_coordinate\_score\_sq\_eq\_hessian}, \texttt{integral\_norm\_sq\_gradient\_eq\_laplacian\_of\_partition}, \texttt{integral\_norm\_sq\_gradient\_eq\_laplacian}.

\section{Differentiate opposite coordinate models to retain the literal M\_j}


Apply the coordinate upper model at x with displacement a e\_j and at x+a e\_j with displacement minus a e\_j. Addition gives a times the difference of the coordinate scores at most M\_j a squared. For a positive, divide by a squared and take the right derivative at zero. Its actual derivative is the diagonal Hessian entry, so partial\_jj f(x) is at most M\_j. Taking the trace proves the pointwise Laplacian bound p M\_av. Conversely strong gradient monotonicity along the same coordinate line bounds its difference quotient below by m. Taking the same derivative gives m at most every diagonal Hessian entry. In the convex case m=0, their sum is nonnegative.



\noindent Lean declarations: \texttt{hessian\_diagonal\_le\_coordinateSmoothness}, \texttt{laplacian\_le\_averageSmoothness}, \texttt{strongConvex\_le\_hessian\_diagonal}, \texttt{laplacian\_nonneg\_of\_convex}.

\section{Use the genuine Laplacian supremum and complete Lemma 7}


Define M\_1Delta as the supremum of the range of the true Laplacian, divided by p. Its range is nonempty, and the pointwise bound p M\_av makes it bounded above. The defining supremum therefore bounds every Laplacian value and is itself at most p M\_av. Dividing by positive p gives M\_1Delta at most M\_av. C2 regularity also makes the Laplacian continuous, as a finite sum of continuous score derivatives. Its nonnegativity and finite pointwise bound prove actual integrability. Under the normalized Gibbs law, the score identity and these pointwise comparisons give exactly

\[\mathbb E_\pi\|\nabla f\|^2=\mathbb E_\pi|\Delta f|\le pM_{1,\Delta}\le pM_{\rm av}.\]

\noindent Lean declarations: \texttt{averageLaplacianSmoothness}, \texttt{averageLaplacianSmoothness\_le\_averageSmoothness}, \texttt{laplacian\_le\_averageLaplacianSmoothness}, \texttt{continuous\_laplacian}, \texttt{score\_laplacian\_identity\_and\_bounds}.

\end{document}
